\section{Limitations}
\label{sec:limits}
첫째, 본 결과가 아직 없다. 이 원고는 가설과 승인된 설계로 쓴 완성 구조이며, 결합 계층의 이득, 학생 상위 모델의 대체 가능성(H-L), 강건성의 주장은 모두 사전 등록한 실험으로 검증할 가설이다. 지금까지의 근거는 한 과제, 작은 표본, 대부분 시뮬의 초기값이다. 둘째, 학습한 8B 상위 모델의 \tentative{4/4}는 학습과 같은 분포의 결과이고 시뮬 참값 라벨로 얻었다. 탁자 높이·물체·방해물이 바뀌면 어떨지, Astra 라벨로도 같은지, 35B에서 더 나은지는 모른다. 셋째, VLA는 아직 단독으로 과제를 끝내지 못한다(\tentative{0/12}). 결정이 실행을 조종하게 하는 레시피도 시뮬에서만 성립했다(실데이터 \tentative{0.536}). 결합의 이득을 재려면 이 바닥부터 넘어야 한다. 넷째, 상위 계획기는 몇 초 전 영상으로 답하고, 7--10\,cm 목표 어긋남을 알아채지 못했다. VLA는 상위가 옳다는 전제로 움직이므로 이 공백이 그대로 전파될 수 있다. 도착 시 대조와 합의는 피해를 줄일 뿐 틀린 계획을 바로잡지 못한다. 다섯째, API 모델은 같은 입력에도 답과 지연이 날마다 달라지고 결정성 설정을 받지 않는다. 여섯째, 목표를 공간에 표현하지 않고 ``어디로''를 수치 명령이나 점에 맡긴 선택은 새 물체에서 점·궤적 조건이 앞섰다는 보고~\cite{steerable2026,molmoact2026}와 긴장한다. 일곱째, 공개 실물 데이터 일부는 라이선스가 명시되지 않아 내부 학습·평가에만 쓰며, 실물 전이는 AI Worker 최소판으로만 확인한다. 여덟째, 안전은 이 논문의 범위가 아니며, 안전 투영과 편향 상한은 결정 권위와 튐 방지를 위한 최소 장치다.

\section{Conclusion}
\label{sec:conclusion}
우리는 느리지만 과제를 아는 상위 계획기와 빠르지만 부정확한 VLA를 잇는 의도 공유 계층 PaceNotes를 설계했다. 랠리의 코드라이버처럼 상위는 다음 구간의 의도를 먼저 선언하고 결과를 검증한다. 드라이버처럼 VLA는 그 의도를 기준으로 실시간 세부 조종과 쥐기·놓기 시점을 맡는다. 느린 답은 도착 시점의 실제 진행과 대조하고, 합의와 거리 권한을 거쳐 서서히 얹는다. 지금까지의 근거는 이 분업의 두 전제를 부분적으로 보여 준다. 상위에 맞춘 인터페이스에서 프런티어 계획기는 혼자서도 조종할 수 있었다. 영샷 오픈급 모델이 실패한 위치 읽기는 무료 시뮬 참값의 촘촘한 감독으로 메워졌다(8B, 분포 안). 동시에 두 공백도 드러났다. 프런티어 계획기는 작은 목표 어긋남을 알아채지 못했고, 지금의 VLA는 단독으로 과제를 끝내지 못한다. 다음 단계는 검증 사다리를 따라 VLA를 다시 학습하고, 결합 형식으로 학습한 학생 상위 모델로 결합 폐루프를 돌리고, 다양화한 데이터와 OOD 세트로 일반성을 재는 것이다. 모든 판정 기준은 실행 전에 고정했다.
